\documentclass[11pt,a4paper]{article}

\usepackage[utf8]{inputenc}
\usepackage[T1]{fontenc}
\usepackage{lmodern}
\usepackage[french]{babel}
\usepackage[margin=1.7cm]{geometry}
\usepackage{microtype}
\usepackage{enumitem}
\usepackage{xcolor}
\usepackage[most]{tcolorbox}

\definecolor{navy}{HTML}{1F3864}
\definecolor{lightblue}{HTML}{D6E0F0}
\definecolor{accent}{HTML}{C0491F}
\definecolor{softgreen}{HTML}{E2EFDA}
\definecolor{softred}{HTML}{FCE4D6}

\newtcolorbox{cle}[1][]{colback=lightblue!45,colframe=navy,boxrule=0.6pt,
  arc=2pt,left=7pt,right=7pt,top=5pt,bottom=5pt,#1}
\newtcolorbox{oui}[1][]{colback=softgreen,colframe=softgreen!60!black,boxrule=0.5pt,
  arc=2pt,left=7pt,right=7pt,top=4pt,bottom=4pt,#1}
\newtcolorbox{non}[1][]{colback=softred,colframe=accent,boxrule=0.5pt,
  arc=2pt,left=7pt,right=7pt,top=4pt,bottom=4pt,#1}

\setlist[itemize]{leftmargin=1.3em,itemsep=1pt,topsep=2pt}
\setlength{\parindent}{0pt}

\title{\vspace{-1.4cm}\color{navy}\bfseries Antisèche d'animation -- séance d'élicitation de $W$\\[2pt]
  \large À garder sous les yeux (~30 min, \textbf{un expert à la fois})}
\date{}

\begin{document}
\maketitle
\vspace{-1.4cm}

\begin{cle}
\textbf{Le principe qui gouverne tout.} On mesure le \textbf{jugement} de l'expert, pas ses
connaissances : il n'y a pas de bonne réponse. Le seul biais qui casse la méthode est
l'\textbf{intervalle trop étroit} (excès de confiance). Ton rôle d'animateur : l'\textbf{élargir}
en poussant sur les extrêmes, \textbf{jamais} en soufflant un chiffre.
\end{cle}

\vspace{4pt}
\textbf{1. Préparer (2 min).}
\begin{itemize}
\item Un expert \textbf{seul}, sans les autres. Aucune discussion entre eux : sinon leurs
      réponses se ressemblent et la pondération de Cooke perd tout son sens.
\item Explique le format sur un exemple \emph{neutre, hors sujet} (ex.\ la hauteur d'un
      immeuble connu) : basse 5\,\%, médiane 50\,\%, haute 95\,\%.
\item Dis que la partie A a des réponses que \emph{tu} connais mais ne révéleras pas, et que
      ça sert à le \textbf{pondérer}, pas à le piéger.
\end{itemize}

\textbf{2. L'ordre des trois nombres, pour chaque question (anti-ancrage).}
\begin{itemize}
\item D'abord la \textbf{HAUTE (95\,\%)} : \og une valeur assez haute pour que tu sois
      surpris, disons 1 chance sur 20, que le vrai soit au-dessus \fg.
\item Puis la \textbf{BASSE (5\,\%)} : pareil, en dessous.
\item La \textbf{MÉDIANE (50\,\%) en DERNIER} : \og une valeur où le vrai a autant de chances
      d'être au-dessus qu'en dessous \fg.
\item \emph{Pourquoi} : demander la médiane d'abord ancre tout autour d'elle. Partir des
      extrêmes \textbf{élargit} les intervalles, c'est exactement le but.
\end{itemize}

\vspace{2pt}
\begin{oui}
\textbf{Relances neutres (à utiliser).}
\begin{itemize}[topsep=1pt]
\item \og Qu'est-ce qui pourrait le faire monter \emph{bien} plus haut ? Et \emph{bien} plus
      bas ? \fg{} (fait imaginer des scénarios extrêmes, casse l'ancrage).
\item \og Serais-tu \emph{vraiment} surpris si le vrai tombait hors de ton intervalle ? \fg{}
      (teste l'honnêteté de l'intervalle ; s'il hésite, élargis).
\end{itemize}
\end{oui}

\begin{non}
\textbf{À bannir absolument.}
\begin{itemize}[topsep=1pt]
\item Donner un chiffre, un ordre de grandeur, ou \og en général les gens disent\dots \fg.
\item \textbf{Réagir} à une réponse : pas de \og ah oui \fg, pas de moue, pas de silence
      appuyé. Le moindre signal l'ancre.
\item Reformuler sa réponse avec une valeur (\og donc environ 50 ? \fg).
\end{itemize}
\end{non}

\vspace{2pt}
\textbf{3. Quand il bloque.} \og Je ne sais pas \fg{} n'est pas une case vide : demande un
intervalle \textbf{large} (\og entre quoi et quoi es-tu à peu près sûr que ça tombe ? \fg).
Un large intervalle honnête est une excellente réponse, pas un échec.

\vspace{3pt}
\textbf{4. Partie B (les liens entre piliers) -- le point le plus sensible.}
\begin{itemize}
\item Rappelle l'échelle à \emph{chaque} lien : $0$ = le pilier $k$ n'entraîne \emph{jamais}
      $j$ ; $1$ = $k$ entraîne \emph{quasi-certainement} $j$.
\item \textbf{Ne suggère jamais de direction} (\og on pense que la gouvernance domine \fg{} est
      interdit). Laisse-le noter P1$\to$P2 et P2$\to$P1 \emph{séparément} :
      \textbf{l'asymétrie qu'on cherche doit venir de lui, pas de toi.}
\end{itemize}

\vspace{3pt}
\textbf{5. Logistique.} Note les trois nombres tout de suite, vérifie
\textbf{basse $<$ médiane $<$ haute}. Environ $1{,}5$ min par question, 16 questions,
${\sim}30$ min. Ne \og corrige \fg{} pas ses réponses après coup : elles sont la donnée.

\vfill
\begin{cle}
\textbf{En une phrase.} Tu es un \emph{miroir neutre} : tu poses, tu relances vers les
extrêmes, tu notes. Dès que tu donnes une info, une valeur ou une réaction, tu fabriques la
réponse au lieu de la mesurer, et le poids de Cooke devient faux.
\end{cle}

\end{document}
